\documentclass{missingbook}

\begin{document}

\frontmatter
\mbtitlepage{Dissolving Problems}{A Casebook of Theory-Building}{Claude}{First edition \textperiodcentered{} October 2026}
\mbtableofcontents

\mainmatter
\casepart{Building the Right Object}{The first family of moves changes \emph{what} we talk about: we complete a space so that limits stop escaping, or we axiomatize a notion so that its hidden parameters must be declared.}{22}

\casechapter{The Needle and the Noodle}{Buffon's problem dissolved by generalization}{Generalize to simplify}{ch:sample}
\begin{refsection}[sample.bib]

\begin{chaptersummary}
Buffon asked for the probability that a needle dropped on ruled paper crosses a line. The classical answer is an integral. Barbier's answer is a sentence: bend the needle into a noodle, and the question becomes so general that it can only have one answer.
\end{chaptersummary}

\section{Before: an integral over angles}

Rule the plane with parallel lines at distance $d$ and drop a needle of length $\ell\le d$ at random.\sidenote{``At random'' means that the distance from the centre of the needle to the nearest line is uniform on $[0,d/2]$, the angle is uniform on $[0,\pi/2]$, and the two are independent.} The classical solution integrates over the angle $\theta$ between the needle and the lines:
\[
  \Pr[\text{cross}] = \frac{2}{\pi}\int_0^{\pi/2}\frac{\ell\sin\theta}{d}\dd\theta = \frac{2\ell}{\pi d}.
\]
The computation is short, but it explains nothing. Why $\pi$? Why is the answer linear in $\ell$? Change the needle to a bent wire and the integral becomes a different integral.\index{Buffon's needle}

\section{The move: let the needle bend}

\begin{definition}[Noodle]\label{def:sample:noodle}
A \term{noodle} is a rigid rectifiable curve of finite length $L$, dropped with a uniformly random position and orientation.
\end{definition}

Let $N(C)$ count the crossings of a noodle $C$ with the lines. The question ``what is $\E N(C)$?'' sounds harder than Buffon's. It is easier, because it has more structure.

\section{After: one line of reasoning}

\begin{theorem}[\textcite{Barbier1860}]\label{thm:sample:barbier}
Every noodle of length $L$ satisfies $\E N(C) = 2L/(\pi d)$.
\end{theorem}

\begin{proof}
The plan is to show that $\E N$ is proportional to length and then to calibrate the constant with a circle. Cut the noodle into pieces; crossings add, so linearity of expectation makes $\E N$ additive over pieces, even though the pieces move together. An additive, monotone function of length is linear: $\E N(C) = cL$. A circle of diameter $d$ meets the lines exactly twice with probability one, and its length is $\pi d$, so $c\cdot\pi d = 2$.
\end{proof}

\begin{corollary}\label{cor:sample:buffon}
A needle of length $\ell\le d$ crosses a line with probability $2\ell/(\pi d)$.
\end{corollary}

\begin{proof}
A needle with $\ell\le d$ crosses at most one line, so $\Pr[\text{cross}] = \E N$, and \cref{thm:sample:barbier} applies.
\end{proof}

\section{Costs and limits}

\paragraph{Where did the difficulty go?} Into linearity of expectation, which holds for dependent pieces, and into the existence of one calibrating shape whose crossing number is constant.

\section{The move, distilled}

\begin{move}{Generalize to simplify}
\movepart{In one sentence} Replace a question about one object by the same question about a family that has more structure; the structure may force the answer.
\movepart{Triggers} A computation that gives the right number but no reason; an answer that is suspiciously linear.
\movepart{Anti-pattern} A generalization whose proof is longer than the special case.
\end{move}

\dissolutiontest{A trigonometric integral becomes a two-line argument.}{Straightness of the needle is no longer used.}{Needles, polygons and curves obey one formula.}{Expected crossings of every rectifiable curve.}

\begin{exercises}
\item \consolidation Prove that an additive monotone function on $[0,\infty)$ is linear.\label{ex:sample:additive}
\item \transfer\hard Use the same move to compute the expected number of crossings of a random line with a convex curve. \hint{Calibrate with a circle again.}
\item \research\harder Which other classical integrals in geometric probability dissolve in the same way?
\end{exercises}

\notesandrefs
The noodle argument is due to \textcite{Barbier1860}; a modern account is in \textcite{AignerZiegler2018}. P\'olya's ``inventor's paradox'' \parencite{Polya1945} is the general principle.

\printchapterbib
\end{refsection}

\printindex
\end{document}
